% !TeX root = ../se200_identity_regimes.tex

\section{Consequences}
\label{sec:consequences}

The result establishes a lower bound of nine on the number of distinct
referential regime profiles required for admissible substrates.
This has several structural consequences.

First, individuation is constrained by invariance:
identity and persistence must be fixed independently of interpretation.
Systems that rely on contextual or role-based reclassification
for identity violate admissibility.

Second, relations operate over entities without determining their identity.
This enforces a separation between reference and interpretation,
ensuring that participation, measurement, and scoping relations
do not alter identity conditions.

Third, regime distinctions are forced by transformation-induced
split pressure. When a transformation must be identity-preserving
under one identity basis and identity-breaking under another,
a single regime is insufficient and refinement is required.
The derived splits of $\mathrm{ENR}$, $\mathrm{CTX}$, and $\mathrm{NOR}$
are instances of this phenomenon.

\subsection{Illustrative Examples}
\label{subsec:examples}

The following examples illustrate the three transformation-induced
splits derived in Sections~\ref{subsec:enr_split}--\ref{subsec:nor_split}.
They serve only to anchor interpretation and do not contribute
to the derivation.

\begin{description}[style=nextline,leftmargin=1.5cm]

  \item[ENR (branching)]
        Consider a physical location and an artifact.
        Under branching, multiple continuations may refer to the same location,
        but distinct artifacts cannot remain identical.
        This forces a distinction between locus-bound and instrument-bound regimes.

  \item[CTX (decomposition)]
        Consider a jurisdiction defined as a union of subregions.
        Decomposition into subcontexts can preserve the same applicability
        while altering internal structure.
        This forces a distinction between extension-based and structure-based contexts.

  \item[NOR (refinement)]
        Consider a normative rule expressed at different levels of detail.
        Refinement can preserve the same normative content
        (what the rule requires or permits)
        while altering its structural organization.
        A content-tracking identity basis classifies such refinement
        as identity-preserving; a structure-tracking basis classifies it
        as identity-breaking.
        This forces a distinction between content-sensitive ($\mathrm{NOR\mbox{-}C}$)
        and structure-sensitive ($\mathrm{NOR\mbox{-}S}$) regimes.

\end{description}

\subsection{Consequences for Representation}
\label{subsec:representation_consequences}

The embedding establishes that all admissible structure is representable
as regime-typed path families over admissible relations.
Accordingly:

\begin{itemize}
  \item structural comparison reduces to comparison of path families,
  \item admissibility reduces to preservation of path structure under transformation,
  \item identity is evaluated at the level of equivalence classes,
        not individual representations.
\end{itemize}

\subsection{Consequences for Interpretation}
\label{subsec:interpretation_consequences}

Interpretive content is not represented at the substrate level.
All interpretive assignments arise in extensions over REC-anchored paths.

Thus:

\begin{itemize}
  \item identical path structure may admit multiple incompatible interpretations,
  \item interpretive disagreement does not affect identity or persistence,
  \item causal and normative content are external to admissible structure.
\end{itemize}

\subsection{Consequences for Comparison}
\label{subsec:comparison_consequences}

Two systems are structurally comparable if and only if
their path families are comparable under regime-typed isomorphism.

Therefore:

\begin{itemize}
  \item comparison is invariant under admissible re-expression and refinement,
  \item differences in interpretation do not affect structural comparison,
  \item non-comparability arises only from differences in regime structure
        or admissible relations.
\end{itemize}

\subsection{Consequences for Intervention}
\label{subsec:intervention_consequences}

Interventions correspond to admissible transformations
of path structure.

An intervention is minimal if it alters the admissibility
or classification of a transformation family
without introducing non-admissible structure.

Accordingly:

\begin{itemize}
  \item modifying $\mathrm{NOR}$ alters admissibility of pathways,
  \item modifying $\mathrm{CTX}$ alters applicability domains,
  \item modifying $\mathrm{ENR}$ alters capacity to instantiate pathways,
  \item modifying $\mathrm{REC}$ alters observability of path structure,
  \item modifying $\mathrm{OCC}$ alters the time-indexed occurrences
      through which path transitions are realized.
\end{itemize}

Interventions propagate along admissible paths and are evaluated
by their effect on path families, not by interpretive outcomes.

\subsection{Consequences for Structural Analysis}
\label{subsec:structural_analysis_consequences}

Structural analysis reduces to the study of path families
under admissible transformations.

In particular:

\begin{itemize}
  \item detection of distortion corresponds to identification of
        anomalous or asymmetric path structure,
  \item explanation corresponds to identification of transformation families
        responsible for such structure,
  \item correction corresponds to minimal modification of regime-level constraints
        affecting those transformation families.
\end{itemize}

No appeal to causal or normative primitives is required
for detection, comparison, or correction at the structural level.
